% from ega2-1-fr.tex
\oldpage[II]{198}
\begin{proposition}[8.14.3]
\label{II.8.14.3-fr}
L'homomorphisme (\hyperref[II.8.14.2.4-fr]{8.14.2.4}) est bijectif~; il en
est de même de (\hyperref[II.8.14.2.3-fr]{8.14.2.3}) lorsque le
$\mathcal{S}$-Module gradué $\mathcal{M}$ admet une présentation finie
(\hyperref[II.3.1.1-fr]{3.1.1}).
\end{proposition}

\begin{proof}
La question est évidemment locale sur $X$ et sur $Y$~; on peut donc supposer
$Y=\operatorname{Spec}(A)$ affine, $\mathcal{S}=\widetilde{S}$,
$\mathcal{M}=\widetilde{M}$, $\mathcal{N}=\widetilde{N}$, où $S$ est une
$A$-algèbre graduée à degrés positifs, $M$ et $N$ deux $S$-modules gradués.
Si $f$ est un élément homogène de $S_+$, les homomorphismes
(\hyperref[II.8.14.2.1-fr]{8.14.2.1}) et
(\hyperref[II.8.14.2.2-fr]{8.14.2.2}), restreints à l'ouvert affine
$D_+(f)$, correspondent aux homomorphismes canoniques
(\hyperref[II.2.5.11.1-fr]{2.5.11.1}) et
(\hyperref[II.2.5.12.1-fr]{2.5.12.1})~:
\begin{align*}
  M(m)_{(f)}\otimes_{S_{(f)}}N(n)_{(f)}
  &\to(M\otimes_S N)(m+n)_{(f)},\\
  (\operatorname{Hom}_S(M,N))(n-m)_{(f)}
  &\to\operatorname{Hom}_{S_{(f)}}(M(m)_{(f)},N(n)_{(f)}).
\end{align*}
Si on se reporte aux définitions de ces homomorphismes, on voit donc (compte
tenu de (\hyperref[II.8.2.9.1-fr]{8.2.9.1})) que la restriction de
(\hyperref[II.8.14.2.4-fr]{8.14.2.4}) à $D_+(f)$ correspond à
l'homomorphisme canonique
\[
  M_f\otimes_{S_f}N_f\to(M\otimes_S N)_f
\]
défini dans (\hyperref[0.1.3.4-fr]{0, 1.3.4}), et on sait que ce dernier est
un isomorphisme. De même, la restriction de
(\hyperref[II.8.14.2.3-fr]{8.14.2.3}) à $D_+(f)$ correspond à
l'homomorphisme canonique (\hyperref[0.1.3.5-fr]{0, 1.3.5})
\[
  (\operatorname{Hom}_S(M,N))_f\to
  \operatorname{Hom}_{S_f}(M_f,N_f)
\]
compte tenu de ce que, $M$ étant de type fini, le module
$\operatorname{Hom}_S(M,N)$, somme directe des sous-groupes formés des
homomorphismes \emph{homogènes} de $S$-modules
(\hyperref[II.2.1.2-fr]{2.1.2}), coïncide avec l'ensemble de \emph{tous} les
homomorphismes $M\to N$ de $S$-modules. L'hypothèse que $M$ admet une
présentation finie entraîne alors (\hyperref[0.1.3.5-fr]{0, 1.3.5}) que
l'homomorphisme canonique considéré est bien un isomorphisme.
\end{proof}

\begin{proposition}[8.14.4]
\label{II.8.14.4-fr}
Si $U$ est un ouvert quasi-compact de $X$, il existe un entier $d$ tel que
pour tout entier $n$ multiple de $d$, $\mathcal{O}_X(n)|U$ soit inversible,
son inverse étant $\mathcal{O}_X(-n)|U$.
\end{proposition}

\begin{proof}
Comme $q(U)$ est quasi-compact, il est recouvert par un nombre fini d'ouverts
affines $V_i$, donc tout $x\in U$ est contenu dans un ouvert affine de la
forme $D_+(f)$, où $f$ est un élément homogène de degré $>0$ de l'un des
anneaux $\Gamma(V_i,\mathcal{S})$. Comme $U$ est quasi-compact, on peut le
recouvrir par un nombre fini de tels ouverts $D_+(f_j)$~; soit $d$ un
multiple commun des degrés des $f_j$. Le nombre $d$ répond à la question en
vertu de (\hyperref[II.2.5.17-fr]{2.5.17}).
\end{proof}

\begin{env}[8.14.5]
\label{II.8.14.5-fr}
Avec les hypothèses et notations de (\hyperref[II.8.14.1-fr]{8.14.1}), on a
défini dans (\hyperref[II.3.3.2-fr]{3.3.2}) des homomorphismes canoniques de
$\mathcal{O}_Y$-Modules
\begin{equation}
\label{II.8.14.5.1-fr}
  \alpha_n:\mathcal{M}_n\to
  q_*(\operatorname{Proj}_0(\mathcal{M}(n)))
  \qquad(n\in\mathbf{Z})
\tag{8.14.5.1}
\end{equation}
Généralisant les notations de (\hyperref[II.3.3.1-fr]{3.3.1}), nous
poserons, pour tout $\mathcal{S}_X$-Module gradué $\mathcal{F}$,
\begin{equation}
\label{II.8.14.5.2-fr}
  \Gamma_*(\mathcal{F})=
  \bigoplus_{n\in\mathbf{Z}}q_*(\mathcal{F}_n).
\tag{8.14.5.2}
\end{equation}
En particulier,
$\Gamma_*(\mathcal{S}_X)=\bigoplus_{n\in\mathbf{Z}}q_*(\mathcal{O}_X(n))$
est la $\mathcal{O}_Y$-Algèbre graduée notée
$\Gamma_*(\mathcal{O}_X)$ dans
(\hyperref[II.3.3.1.2-fr]{3.3.1.2})~; il est clair que
$\Gamma_*(\mathcal{F})$ est un $\Gamma_*(\mathcal{S}_X)$-Module gradué
(\hyperref[0.4.2.2-fr]{0, 4.2.2}). Lorsque

